\section{Isolated equations for the twenty curved routes}\label{sec:roots}
Write $p,q$ for the two developed points in $(U,V)$ order, $T=2\pi/3$
and $G(f)$ for a gate with $0\le f\le1$. A curved gate uses its own
normalized arclength, so $G_A(f)=a[Tf]$ and
$G_B(f)=(W,H)-a[Tf]$. For a straight gate $L$ or $R$ let $J$
be reflection in its entire supporting line. The following equations
specify exact values independently of a choice of optimizer.

For each winning word containing just one curved gate, take $P=p$ and $Q=q$.
If the word starts with a straight gate, replace $P$ by $Jp$.
If it ends with a straight gate, replace $Q$ by $Jq$. Write $C$ for the
curved gate in the word. There is precisely one $f$ in its interval in
Table~\ref{tab:rootboxes} satisfying
\begin{equation}
 C'(f)\cdot\left(\frac{C(f)-P}{|C(f)-P|}
             +\frac{C(f)-Q}{|C(f)-Q|}\right)=0 .
 \label{eq:scalarcurve}
\end{equation}
The distance equals $|P-C(f)|+|C(f)-Q|$.
Indeed, when a straight gate occurs, its certified minimizing contact is
interior. For the curved contact held fixed, differentiating in the
straight contact shows that its two unit directions have equal tangential
components. Equivalently, reflection puts the two legs on one straight
line. One can also see this by minimizing the sum with the contact on
the entire supporting line. The curved contact and the other endpoint
are strictly on the same side, so the segment to the reflected endpoint
meets the line in the unique unconstrained minimizer. Its distance sum
is strictly convex as a function of the line contact, so the interior
constrained minimizer is that intersection.

For $P_4,P_{10}$ put $x=G_A(f),y=G_B(g)$. The two equations are
\begin{align}
 G_A'(f)\cdot\left(\frac{x-p}{|x-p|}+\frac{x-y}{|x-y|}\right)&=0,\nonumber\\
 G_B'(g)\cdot\left(\frac{y-x}{|y-x|}+\frac{y-q}{|y-q|}\right)&=0.
 \label{eq:doublecurve}
\end{align}
There is precisely one solution in the listed box and the distance is
$|p-x|+|x-y|+|y-q|$.

\begin{table}[!ht]\centering\small
\caption{Isolating intervals in gate order. An integer $k$ denotes the
closed interval $[k,k+3]/4096$. For the scalar equation use the integer
in the curved gate's position.}\label{tab:rootboxes}
\begin{tabular}{@{}ccl@{\qquad}ccl@{}}\toprule
Pair & Word & Interval indices & Pair & Word & Interval indices\\\midrule
$1,4$ & A & 1317 & $3,10$ & RB & 1233, 3409\\
$1,7$ & AL & 834, 2515 & $4,5$ & LA & 2119, 3575\\
$1,8$ & AL & 688, 1536 & $4,6$ & AL & 2600, 3814\\
$2,4$ & A & 3123 & $4,9$ & A & 3259\\
$2,6$ & B & 4052 & $4,10$ & AB & 3015, 3521\\
$2,7$ & BR & 3871, 1561 & $5,7$ & A & 4060\\
$2,10$ & B & 3498 & $5,8$ & A & 3625\\
$3,5$ & LA & 3853, 3766 & $6,8$ & A & 4070\\
$3,6$ & RB & 1075, 4020 & $7,9$ & BR & 2839, 1172\\
$3,7$ & B & 2981 & $9,10$ & B & 3488\\\bottomrule
\end{tabular}
\end{table}

\paragraph{Why the equations isolate the global shortest paths.}
Every minimizing tuple of the winning word lies in its box. To prove this
finite assertion, \texttt{verify\_roots.py} replays the dyadic covers in
\texttt{root\_trees.json} using the same chord-tube support bounds as the
distance certificate. Each leaf is contained in the displayed box or has
a lower bound at least one unit of $10^{-14}$ above the feasible path
upper bound. There are $1{,}457$ such leaves across the twenty covers.
All boxes are interior to the admissible cubes, so a minimizing tuple
satisfies the ordinary stationarity equations. The following positive
Hessian certificate proves uniqueness on the stated box.

For a leg displacement $r=z_{j+1}-z_j$ let $r_i$ be its derivative with
respect to contact coordinate $i$, and $r_{ih}$ its mixed derivative.
The contribution to the Hessian is
\[
 \partial_{ih}|r|=
 \frac{(r\times r_i)(r\times r_h)}{|r|^3}
       +\frac{r\cdot r_{ih}}{|r|},
 \qquad (a,b)\times(c,d)=ad-bc .
\]
All mixed $r_{ih}$ with $i\ne h$ vanish. The nonzero $r_i$ and $r_{ii}$
are the signed first and second derivatives of the incident gate.
For $A$ put $v=\tan(Tf/2)/\sqrt3$, $w=\sqrt{1-v^2}$,
$d=1+3v^2$, $\tau=(w^3,v(3+v^2))/d$. With $J_0(a,b)=(-b,a)$ one has
\[
 G_A'=T\tau,\qquad G_A''=T^2(\sqrt3 w/2)J_0\tau .
\]
The derivatives of $B$ are their negatives, both straight gates have
derivative $(0,\sqrt3)$ and their second derivatives vanish.
To bound positions on an interval it suffices to enclose their two
endpoints coordinatewise because each coordinate is monotone. The
identity $v^2=(1-\cos(Tf))/(3(1+\cos(Tf)))$ encloses all derivative
terms by interval operations. On every stated box each original and
unfolded leg has length greater than $0.064$. Summing the formula
above with outward rounding gives $D^2 L\succ I/32$. In dimension two
the checker proves positivity of the first principal minor and determinant
of $D^2L-I/32$. It separately applies the same check to the one-variable
reflected expression in each mixed case on its entire curved interval.
Integrating the Hessian on the segment joining any two distinct stationary
points would give
$(\nabla L(b)-\nabla L(a))\cdot(b-a)\ge |b-a|^2/32>0$.
Thus two such points cannot exist. This argument also supplies the
quantitative estimate
\[
       0\le L(u)-L(u_*)\le 16\,|\nabla L(u)|^2
\]
for every $u$ in the original box and its minimizing contact $u_*$. It
follows by integrating the Hessian lower bound from $u$ to $u_*$ and
completing the square.
